% CP-VI-0032
% target_chapter: VI/18 — Affine Schemes
% source_unit_id: IMP-DL-D27460A55B-P001
% source: Downloads/theory-of-algebraic-geometry-5.tex L193-L277
% solution_provenance: SOURCE_EXPLICIT_SOLUTION

\subsection*{Problem 2 - Intersections of Affine Open Subschemes}

\textbf{Problem.}
Let \(X\) be a scheme, with open affine subsets
\[
U=\operatorname{Spec} A,
\qquad
V=\operatorname{Spec} B.
\]
Show that \(U\cap V\) can be covered by open affine subschemes \(\{U_i\}\) such that there are elements
\[
f_i\in A,
\qquad
g_i\in B,
\]
with
\[
U_i=D(f_i)\subseteq U
\qquad\text{and}\qquad
U_i=D(g_i)\subseteq V.
\]

\textbf{Solution.}

We use the fact that distinguished open subsets form a basis for the topology of an affine scheme.

Let
\[
x\in U\cap V.
\]
Since \(U\cap V\) is open in \(U=\operatorname{Spec} A\), and distinguished opens form a basis for \(\operatorname{Spec} A\), there exists \(f\in A\) such that
\[
x\in D(f)\subseteq U\cap V.
\]
Thus \(D(f)\) is an affine open subset of \(U\), and hence an affine open subscheme of \(X\).

Now \(D(f)\subseteq V=\operatorname{Spec} B\). Since \(D(f)\) is open in \(V\), and distinguished opens form a basis for \(V\), there exists \(g\in B\) such that
\[
x\in D(g)\subseteq D(f)\subseteq V.
\]
Here \(D(g)\) is a distinguished open subset of \(V\).

We claim that the chosen open \(D_V(g)\) is also distinguished in
\(U\), without a second arbitrary shrinking. Since
\(D_U(f)\subseteq V\), the regular function \(g\in B=\Gamma(V,\mathcal O_V)\)
restricts to a section on \(D_U(f)\). Under
\[
\Gamma(D_U(f),\mathcal O_X)=A_f,
\]
write this restriction as \(g|_{D_U(f)}=h/f^N\) with \(h\in A\) and
\(N\geq0\). At every point of \(D_U(f)\), the germ of \(f\) is a unit;
therefore the locus where \(g\) is invertible is exactly the locus
where \(h\) is invertible. In other words,
\[
D_V(g)\cap D_U(f)=D_U(fh).
\]
Because we chose \(D_V(g)\subseteq D_U(f)\), this gives the
\emph{equality of open subsets of \(X\)}
\[
D_V(g)=D_U(fh).
\]
Thus \(U_x:=D_V(g)=D_U(fh)\) is a neighborhood of \(x\),
distinguished in both \(V\) and \(U\). This establishes the
simultaneous distinguished-open condition rather than merely finding
a smaller distinguished open on one side.

As \(x\) varies over \(U\cap V\), the sets \(U_x\) cover \(U\cap V\). Hence \(U\cap V\) can be covered by affine opens \(\{U_i\}\) such that
\[
U_i=D(f_i)\subseteq U,
\qquad
U_i=D(g_i)\subseteq V.
\]

Thus
\[
\boxed{
	U\cap V
	\text{ can be covered by affine opens distinguished in both }U\text{ and }V.
}
\]
